\section{Retained matrices and finite families of ordered cores}

The identities below give supporting algebra for
\cite[Proposition~5.7]{Arad2016RigorousRG}.
All products of site matrices retain their natural order.

\begin{definition}[An interval and its single retained complement]
    \label{def:retained_interval_coordinates}
    \lean{TNLean.MPS.AreaLaw.retainedIntervalEquiv,MPOTensor.actualAffineRetainedConfig}
    Let $T,L,a$ be nonnegative integers with $a+L\le T$, and set
    $J=[a,a+L)$ in the chain of $T$ sites. For any set $X$, restriction gives
    an equivalence
    \begin{align}
        e:X^{J^c}\times X^{\{0,\ldots,L-1\}} & \longrightarrow
        X^{\{0,\ldots,T-1\}},\notag                            \\
        e(\eta,x)_v                          & =
        \begin{cases}x_{v-a},&a\le v<a+L,\\ \eta_v,&v\notin J.\end{cases}
        \label{eq:retained_interval_coordinates}
    \end{align}
    Its inverse restricts a configuration to $J^c$ and to the sites $a+r$,
    $0\le r<L$. For local alphabet $\{0,\ldots,q-1\}$, denote the retained
    configuration set by $\mathcal R=\{0,\ldots,q-1\}^{J^c}$.
    The complement is a single set: each site outside $J$ occurs once.
\end{definition}

\begin{definition}[The actual retained closing matrix]
    \label{def:actual_retained_interval_closing}
    \lean{MPOTensor.actualRetainedIntervalClosing}
    \uses{def:retained_interval_coordinates}
    Let $A_v(i,j)\in\C^{B\times B}$ be supplied site-dependent operator
    tensors on the full chain, and let $C\in\C^{B\times B}$ be supplied.
    For $\eta,\theta\in\mathcal R$, define
    \begin{align}
        P(\eta,\theta) & =\prod_{v=0}^{a-1}A_v(\eta_v,\theta_v),         &
        S(\eta,\theta) & =\prod_{v=a+L}^{T-1}A_v(\eta_v,\theta_v),\notag   \\
        F(\eta,\theta) & =S(\eta,\theta)\,C\,P(\eta,\theta).
        \label{eq:actual_retained_interval_closing}
    \end{align}
    Empty products are identity matrices. The matrix $F(\eta,\theta)$ is
    chosen before any physical coordinates inside $J$.
\end{definition}


\begin{theorem}[Restriction of a trace contraction to an interval]
    \label{thm:actual_retained_interval_trace}
    \lean{MPOTensor.trace_actual_retainedInterval_eq}
    \uses{def:retained_interval_coordinates,def:actual_retained_interval_closing,
        def:trace_chain_operator}
    Use $a+L\le T$, the full-chain tensors $A_v$, the closing matrix $C$,
    the coordinate equivalence $e$, and the fixed retained configurations
    $\eta,\theta$ of (\ref{eq:actual_retained_interval_closing}).
    Put $A^{\mathrm{int}}_r=A_{a+r}$ for $0\le r<L$.
    For all interior configurations $\sigma,\tau$,
    \begin{align}
        \mathcal T(A,C)(e(\eta,\sigma),e(\theta,\tau))
        &=\mathcal T(A^{\mathrm{int}},F(\eta,\theta))(\sigma,\tau).
        \label{eq:actual_retained_interval_trace}
    \end{align}
    Empty intervals and zero dimensions are allowed whenever the supplied
    configurations exist.
\end{theorem}
\begin{proof}
    Split the ordered product into prefix, interval and suffix.
    Cyclicity gives $\operatorname{tr}(CPXS)=\operatorname{tr}(SCPX)$,
    and $F=SCP$ by (\ref{eq:actual_retained_interval_closing}).
\end{proof}

\begin{theorem}[Insertion of fixed configurations at the same trace dimension]
    \label{thm:retained_single_trace_insert}
    \lean{MPOTensor.exists_traceChainOperator_retainedSingle_insert}
    \uses{def:retained_interval_coordinates,def:trace_chain_operator}
    Let $T,L,q,\chi,a$ be nonnegative integers with $a+L\le T$.
    Put $J=[a,a+L)$ and supply site tensors $A_r$ on the $L$ interior sites,
    of virtual dimension $\chi$, and a closing matrix $C$.
    Fix row and column configurations $\eta,\theta$ on $J^c$.
    There are tensors $A'_v$ on the full $T$ sites, of the same dimension,
    such that for all full-chain configurations $\sigma,\tau$,
    \begin{align}
        \mathcal T(A',C)(\sigma,\tau)
        &=\mathbf1_{\sigma|_{J^c}=\eta}\mathbf1_{\tau|_{J^c}=\theta}
        \mathcal T(A,C)(\sigma|_J,\tau|_J).
        \label{eq:retained_single_trace_insert}
    \end{align}
    The restrictions to $J$ use its increasing enumeration. Both $L$ and
    $T$ may vanish; no positivity assumption on $q$ or $\chi$ is imposed.
\end{theorem}
\begin{proof}
    Use $A'_v=A_{v-a}$ inside $J$ and
    $A'_v(i,j)=\mathbf1_{i=\eta_v}\mathbf1_{j=\theta_v}\Id_\chi$ outside $J$.
    If the retained configurations match, the extra letters are identity
    matrices and the trace is unchanged. Otherwise one letter is zero.
\end{proof}

\begin{theorem}[A trace contraction after restriction and fixed extension]
    \label{thm:trace_interval_slice_configuration_extend}
    \lean{MPOTensor.exists_traceChainOperator_intervalSlice_extend_by_configuration}
    \uses{thm:actual_retained_interval_trace,thm:retained_single_trace_insert}
    Let $T,L,N,q,\chi,a,b$ be nonnegative integers with $a+L\le T$ and
    $b+L\le N$. Supply old tensors $A_v$ of dimension $\chi$, a closing
    matrix $C$, and full old configurations $f,g$. For an interior
    configuration $z$, write $\iota_a(f,z)$ for the old configuration
    equal to $z_{v-a}$ on $[a,a+L)$ and to $f_v$ elsewhere.
    Fix new outside configurations $\eta,\theta$ on
    $J'^c$, where $J'=[b,b+L)$ in the new chain of $N$ sites.
    There are a matrix $C'$ and tensors $A'_v$ of the same virtual
    dimension $\chi$ such that
    \begin{align}
        \mathcal T(A',C')(\sigma,\tau)
        &=\mathbf1_{\sigma|_{J'^c}=\eta}\mathbf1_{\tau|_{J'^c}=\theta}
        \mathcal T(A,C)\bigl(\iota_a(f,\sigma|_{J'}),
            \iota_a(g,\tau|_{J'})\bigr).
        \label{eq:trace_interval_slice_configuration_extend}
    \end{align}
    Restrictions to $J'$ use increasing coordinates. All chain lengths and
    dimensions may vanish, subject only to the stated supplied data.
\end{theorem}
\begin{proof}
    \uses{thm:actual_retained_interval_trace,thm:retained_single_trace_insert}
    Fix the old outside configurations to those of $f,g$ and apply
    Theorem~\ref{thm:actual_retained_interval_trace}, obtaining the restricted
    tensors and their suffix-closing-prefix matrix $C'$.
    Insert the new retained configurations by
    Theorem~\ref{thm:retained_single_trace_insert}, retaining $C'$ and $\chi$.
\end{proof}

\begin{definition}[The scalar-closed core at a retained matrix entry]
    \label{def:retained_stack_core}
    \lean{MPOTensor.retainedStackCore,MPOTensor.exists_retainedStackCore_traceChainOperator}
    \uses{def:core_word_coefficients,def:trace_core_word_closing}
    Let $m$ be a finite retained index set, let
    $R_{ab}\in\C^{m\times m}$ for $0\le a,b<B$, and let
    $s=((K_0,c_0),\ldots,(K_{h-1},c_{h-1}))$ be a list of affine steps.
    For $t\ge0$ and $i,j\in m$, let
    \begin{align}
        \gamma_t^{i,j}(w)=C_R(s,w)_{ij},
        \qquad w\in(\{0,\ldots,B-1\}^2)^t.
        \label{eq:retained_stack_scalar_coefficient}
    \end{align}
    The joint closing matrix $D_t^{i,j}$ is the closing matrix associated to
    these scalar word coefficients. In the stacking coordinates $\iota_t$
    on words of $t$ virtual labels, its entries are
    \begin{align}
        (D_t^{i,j})_{\iota_t(v),\iota_t(u)}
        =\gamma_t^{i,j}\bigl((u_r,v_r)_{r=0}^{t-1}\bigr).
        \label{eq:retained_stack_closing_entries}
    \end{align}
    For supplied interior tensors $A_v$, $0\le v<L$, let $A_v^{[t]}$ be
    their existing $t$-layer stack, of virtual dimension $B^t$. Define
    \begin{align}
        W_t^{i,j}(\sigma,\tau)=\operatorname{tr}\left(
        D_t^{i,j}\prod_{v=0}^{L-1}A_v^{[t]}(\sigma_v,\tau_v)\right).
        \label{eq:retained_stack_core}
    \end{align}
    The closing matrix depends on $R,s,t,i,j$, independently of the interior
    chain. In particular, no rank-one condition on that matrix is imposed.
\end{definition}

\begin{theorem}[A finite family in natural physical order]
    \label{thm:finite_natural_retained_core_family}
    \lean{MPOTensor.exists_finite_natural_retained_core_family}
    \uses{def:retained_interval_coordinates,def:retained_stack_core,
        def:trace_core_matrix_slice,thm:core_word_product,
        thm:trace_core_slice_words,thm:trace_core_open_square,
        thm:chain_fixed_configuration_extension}
    Fix $a+L\le T$, interior tensors $A_v$ of virtual dimension $B$, and
    a retained index set $m=\mathcal R$. Choose $R_{ab}\in\C^{m\times m}$
    and a list $s$ of $h$ affine steps on $m$.
    Let $N_{ab}(A)$ be the interior matrix slice obtained by fixing the two
    virtual endpoints to $a,b$. Put $n=(h+1)|m|^2$.
    There are degrees $t_r\in\{0,\ldots,h\}$, retained indices $i_r,j_r\in m$,
    and full-chain operators $V_r$, for $0\le r<n$, such that
    \begin{align}
        V_r & =\left(|i_r\rangle\langle j_r|\otimes
        W_{t_r}^{i_r,j_r}\right)^e,\notag           \\
        \left(\prod_{v=0}^{h-1}\left(K_v\otimes I+
        \sum_{a,b}(c_vR_{ab})\otimes N_{ab}(A)\right)\right)^e
            & =\sum_{r=0}^{n-1}V_r.
        \label{eq:finite_natural_retained_core_family}
    \end{align}
    The superscript denotes transport from retained/interior coordinates by
    the equivalence $e$. Each $V_r$, viewed as a vector on the paired physical
    alphabet, has an open-boundary MPS representation with bond bound
    $(B^{t_r})^2$. Empty chains, empty intervals, and zero physical or virtual
    dimensions are included whenever the stated retained data exist.
\end{theorem}
\begin{proof}
    \uses{def:retained_interval_coordinates,def:retained_stack_core,
        thm:core_word_product,thm:trace_core_slice_words,thm:trace_core_open_square,
        thm:chain_fixed_configuration_extension}
    Expand the ordered affine product over words of each length $0\le t\le h$.
    Resolve every retained coefficient on the matrix units of $m$. For fixed
    $t,i,j$, the scalar word coefficients are exactly
    (\ref{eq:retained_stack_scalar_coefficient}). Their joint closing matrix
    therefore gives the core (\ref{eq:retained_stack_core}). Enumerate the
    finite set $\{0,\ldots,h\}\times m\times m$ to obtain the literal family
    in (\ref{eq:finite_natural_retained_core_family}).

    A trace-closed chain of virtual dimension $B^t$ has open-boundary bond
    bound $(B^t)^2$. Inserting fixed retained configurations outside the
    interval gives bond bound $\max\{1,(B^t)^2\}$. If $B^t>0$, this is the
    stated square bound. If $B^t=0$, the trace defining the core vanishes;
    the full inserted operator is zero and has the square bound as well.
\end{proof}

\begin{theorem}[The affine product of the actual supplied trace operator]
    \label{thm:actual_natural_affine_product}
    \lean{MPOTensor.prod_actual_natural_affine_eq}
    \uses{def:retained_interval_coordinates,def:actual_retained_interval_closing}
    Fix $a+L\le T$. For supplied full-chain tensors $A^{\mathrm{full}}_v$
    of virtual dimension $B$ and cap $C\in\C^{B\times B}$, define
    \begin{align}
        H(\sigma,\tau)      & =\operatorname{tr}\left(
        C\prod_{v=0}^{T-1}A^{\mathrm{full}}_v(\sigma_v,\tau_v)\right),\notag \\
        A^{\mathrm{int}}_r  & =A^{\mathrm{full}}_{a+r}, &
        R_{ab}(\eta,\theta) & =F(\eta,\theta)_{ba}.
        \label{eq:actual_natural_affine_data}
    \end{align}
    For every ordered list $s=((K_v,c_v))_{v=0}^{h-1}$ on $\mathcal R$,
    \begin{align}
        \prod_{v=0}^{h-1}\left((K_v\otimes I)^e+c_vH\right)
        =\left(\prod_{v=0}^{h-1}\left(K_v\otimes I+
        \sum_{a,b}(c_vR_{ab})\otimes N_{ab}(A^{\mathrm{int}})\right)\right)^e.
        \label{eq:actual_natural_affine_product}
    \end{align}
    Here $F$ is exactly the suffix-cap-prefix matrix of
    (\ref{eq:actual_retained_interval_closing}). The identity concerns the
    supplied trace operator $H$; no identification with another Hamiltonian
    is assumed.
\end{theorem}
\begin{proof}
    \uses{def:retained_interval_coordinates,def:actual_retained_interval_closing}
    Split the ordered full-chain product into prefix, interval and suffix.
    Cyclicity of the trace moves the suffix before the cap, giving $F=SCP$.
    Expanding the trace over its two virtual indices gives
    $H^{{e}^{-1}}=\sum_{a,b}R_{ab}\otimes N_{ab}(A^{\mathrm{int}})$;
    the reversed indices in $F_{ba}$ come from the trace pairing.
    Transport preserves sums, scalar multiplication and ordered products.
    Applying it to each affine factor proves
    (\ref{eq:actual_natural_affine_product}).
\end{proof}
